\section{סיכום הרצאה: חיבור תנע זוויתי ומקדמי קלבש-גורדן}
\subsection{מבוא: הבעיה של חיבור תנע זוויתי}
במכניקה קלאסית, כאשר יש לנו שני חלקיקים עם תנע זוויתי $\vec{J}_1$ ו-$\vec{J}_2$, התנע הזוויתי הכולל הוא פשוט סכום וקטורי: $\vec{J} = \vec{J}_1 + \vec{J}_2$.
במכניקת הקוונטים, המצב מורכב יותר מכיוון שהאופרטורים אינם בהכרח מתחלפים. אנו מעוניינים למצוא את המצבים העצמיים של התנע הזוויתי הכולל של המערכת.

נניח שני חלקיקים (או שתי דרגות חופש, כגון ספין ותנע זוויתי אורביטלי) הנמצאים במצבים המוגדרים על ידי:
\begin{itemize}
    \item חלקיק 1: תנע זוויתי $j_1$ והיטל $m_1$.
    \item חלקיק 2: תנע זוויתי $j_2$ והיטל $m_2$.
\end{itemize}

\section{הבסיסים השונים}

\subsection{הבסיס הלא-מצומד (\LR{The Uncoupled Basis})}
המרחב של המערכת הוא מכפלה טנזורית של המרחבים של כל חלקיק. הבסיס הטבעי הראשוני הוא בסיס המכפלה:
\[
\ket{j_1 m_1; j_2 m_2} \equiv \ket{j_1 m_1} \otimes \ket{j_2 m_2}
\]
בסיס זה הוא אוסף המצבים העצמיים המשותפים של ארבעה אופרטורים מתחלפים (בהנחה ש-$j_1, j_2$ קבועים):
\[
\vb{J}_1^2, \quad \vb{J}_{1z}, \quad \vb{J}_2^2, \quad \vb{J}_{2z}
\]
הממדים של המרחב הם $(2j_1+1)(2j_2+1)$.

\subsection{הבסיס המצומד (\LR{The Coupled Basis})}
אנו מגדירים את אופרטור התנע הזוויתי הכולל:
\[
\vb{J} = \vb{J}_1 + \vb{J}_2
\]
אנו מעוניינים בבסיס המאופיין על ידי גודל התנע הזוויתי הכולל $J$ וההיטל שלו $M$. בסיס זה, המסומן כ-$\ket{j_1 j_2; J M}$ (או בקיצור $\ket{J M}$), הוא בסיס של וקטורים עצמיים של האופרטורים:
\[
\vb{J}^2, \quad \vb{J}_z, \quad \vb{J}_1^2, \quad \vb{J}_2^2
\]

\subsubsection{מדוע $\vb{J}^2$ אינו אלכסוני בבסיס הלא-מצומד?}
נבחן את האופרטור $\vb{J}^2$:
\[
\vb{J}^2 = (\vb{J}_1 + \vb{J}_2)^2 = \vb{J}_1^2 + \vb{J}_2^2 + 2\vb{J}_1 \cdot \vb{J}_2
\]
האיבר הבעייתי הוא האיבר המעורב (איבר האינטראקציה) $2\vb{J}_1 \cdot \vb{J}_2$, שניתן לכתוב אותו כ:
\[
2\vb{J}_1 \cdot \vb{J}_2 = 2J_{1z}J_{2z} + J_{1+}J_{2-} + J_{1-}J_{2+}
\]
האופרטורים $J_{1+}$ ו-$J_{2-}$ משנים את ערכי ה-$m_1$ ו-$m_2$. לכן, המצבים $\ket{j_1 m_1; j_2 m_2}$ אינם מצבים עצמיים של $\vb{J}^2$ (אלא אם כן מדובר במצבי הקצה), ונדרש מעבר בסיס (לכסון) כדי למצוא את ה-$\ket{J M}$.

\section{מקדמי קלבש-גורדן \LR{Clebsch-Gordan Coefficients}}

המעבר בין הבסיס הלא-מצומד לבסיס המצומד הוא טרנספורמציה אוניטרית. כיוון ששני הבסיסים פורשים את אותו המרחב, ניתן לכתוב את המצבים המצומדים כסופרפוזיציה של המצבים הלא-מצומדים:
\[
\ket{J M} = \sum_{m_1, m_2} \ket{j_1 m_1; j_2 m_2} \braket{j_1 m_1; j_2 m_2}{J M}
\]
המקדמים $C(j_1 m_1 j_2 m_2; J M) \equiv \braket{j_1 m_1; j_2 m_2}{J M}$ נקראים \textbf{מקדמי קלבש-גורדן}.
נהוג לבחור את הפאזות כך שהמקדמים יהיו \textbf{ממשיים}.

\subsection{תכונות אורתוגונליות וסגירות}
כיוון שמדובר במטריצה אורתוגונלית (ממשית ואוניטרית), מתקיימים יחסי השלמות והאורתוגונליות:
\begin{enumerate}
    \item \textbf{אורתוגונליות (סכום על $m_1, m_2$):}
    \[
    \sum_{m_1, m_2} \braket{j_1 m_1 j_2 m_2}{J M} \braket{j_1 m_1 j_2 m_2}{J' M'} = \delta_{JJ'} \delta_{MM'}
    \]
    \item \textbf{סגירות (סכום על $J, M$):}
    \[
    \sum_{J, M} \braket{j_1 m_1 j_2 m_2}{J M} \braket{j_1 m_1' j_2 m_2'}{J M} = \delta_{m_1 m_1'} \delta_{m_2 m_2'}
    \]
\end{enumerate}

\section{כללי ברירה וטווח הערכים האפשריים}

\subsection{כלל ראשון: חיבור ההיטלים}
האופרטור $J_z = J_{1z} + J_{2z}$ הוא אלכסוני בשני הבסיסים. הפעלתו על שני צדי משוואת ההגדרה מראה כי המקדם מתאפס אלא אם כן:
\[
M = m_1 + m_2
\]

\subsection{כלל שני: אי-שוויון המשולש (טווח ה-$J$)}
אילו ערכים של $J$ יכולים להתקבל מחיבור $j_1$ ו-$j_2$?
נשתמש בטיעון של ספירת מצבים (\LR{Degeneracy Argument}).
נסמן ב-$n(M)$ את מספר המצבים הבלתי-תלויים בבסיס המכפלה שיש להם היטל כולל $M$.
נסמן ב-$N(J)$ את מספר הפעמים שהצגה עם תנע זוויתי $J$ מופיעה בפירוק (הניוון של ה-$J$).

אנו יודעים שלכל מולטיפלט עם תנע $J$ יש בדיוק מצב אחד עם היטל $M$, לכל $M$ בטווח $-J \le M \le J$.
לכן, סך המצבים עם היטל $M$ נתון הוא סכום המצבים מכל ה-$J$-ים ש"גדולים מספיק" כדי להכיל את ה-$M$ הזה:
\[
n(M) = \sum_{J \ge |M|} N(J)
\]
מתוך כך נובע כי מספר המולטיפלטים עם ערך ספציפי של $J$ הוא:
\[
N(J) = n(M=J) - n(M=J+1)
\]
\textbf{ניתוח דיאגרמה (המלבן):}
אם נצייר דיאגרמה של כל הזוגות $(m_1, m_2)$, נראה כי:
\begin{enumerate}
    \item הערך המקסימלי האפשרי ל-$M$ הוא $m_1^{max} + m_2^{max} = j_1 + j_2$. יש רק זוג אחד כזה ($(j_1, j_2)$). לכן $n(j_1+j_2)=1$. זה מחייב קיום של $J_{max} = j_1+j_2$ וש-$N(J_{max})=1$.
    \item ככל שמקטינים את $M$ ביחידה אחת, הניוון $n(M)$ גדל ב-1, כל עוד אנו לא "יוצאים" מהתחום (כלומר עד $|j_1 - j_2|$).
    \item חישוב ההפרש $n(J) - n(J+1)$ מראה כי לכל $J$ בתחום המותר, $N(J)=1$. כלומר, כל $J$ מופיע בדיוק פעם אחת.
\end{enumerate}

\textbf{המסקנה - הטווח המותר:}
\[
|j_1 - j_2| \le J \le j_1 + j_2
\]
כמו כן, הסכום $j_1 + j_2 + J$ חייב להיות מספר שלם.

\section{דוגמה: חיבור שני חלקיקים עם ספין 1/2}
נתון: $j_1 = 1/2, \quad j_2 = 1/2$.
הערכים האפשריים ל-$J$ הם לפי אי-שוויון המשולש:
\[
1/2 - 1/2 \le J \le 1/2 + 1/2 \quad \Rightarrow \quad J = 0, 1
\]
\textbf{המצבים:}
\begin{itemize}
    \item \textbf{טריפלט ($J=1$):} מצבים סימטריים.
    \begin{align*}
        \ket{1, 1} &= \ket{\uparrow \uparrow} \quad (\text{כי זה המצב עם ה-M המקסימלי}) \\
        \ket{1, 0} &= \frac{1}{\sqrt{2}}(\ket{\uparrow \downarrow} + \ket{\downarrow \uparrow}) \quad (\text{על ידי הפעלת } J_-) \\
        \ket{1, -1} &= \ket{\downarrow \downarrow}
    \end{align*}
    \item \textbf{סינגלט ($J=0$):} מצב אנטי-סימטרי.
    \begin{align*}
        \ket{0, 0} &= \frac{1}{\sqrt{2}}(\ket{\uparrow \downarrow} - \ket{\downarrow \uparrow}) \quad (\text{מאורתוגונליות ל-} \ket{1,0})
    \end{align*}
\end{itemize}

\section{סימונים נוספים: סמלי ויגנר \LR{3-j}}לעיתים קרובות משתמשים בסימון אלטרנטיבי בעל תכונות סימטריה יפות יותר, הנקרא \textbf{סימני 3-j של ויגנר}. הקשר למקדמי קלבש-גורדן הוא:
\[
\begin{pmatrix}
j_1 & j_2 & j_3 \\
m_1 & m_2 & m_3
\end{pmatrix}
= \frac{(-1)^{j_1 - j_2 - m_3}}{\sqrt{2j_3 + 1}} \braket{j_1 m_1 j_2 m_2}{j_3, -m_3}
\]
שימו לב שהסימון כולל היפוך סימן ב-$m_3$ והתנאי הוא $m_1+m_2+m_3=0$.
\textbf{תכונות סימטריה:}
\begin{itemize}
    \item אינווריאנטיות תחת הזזה ציקלית של העמודות.
    \item כפל בפאזה $(-1)^{j_1+j_2+j_3}$ תחת החלפה של שתי עמודות.
    \item כפל בפאזה תחת היפוך כל הסימנים של ה-$m$-ים.
\end{itemize}

\section{חיבור של שלושה תנעים זוויתיים ומקדמי רקח}
כאשר מחברים שלושה תנעים ($j_1, j_2, j_3$), סדר החיבור משנה את בסיס הביניים אך לא את התוצאה הפיזיקלית הסופית (המרחב הנפרש זהה).
\begin{enumerate}
    \item מסלול א': לחבר $j_1+j_2 \to J_{12}$, ואז לחבר $J_{12}+j_3 \to J$.
    \item מסלול ב': לחבר $j_2+j_3 \to J_{23}$, ואז לחבר $j_1+J_{23} \to J$.
\end{enumerate}
המקדמים המקשרים בין שני הבסיסים הללו נקראים \textbf{מקדמי רקח (\LR{Racah coefficients})} או \textbf{סימני 6-j של ויגנר}. הם תלויים ב-6 מספרי תנע זוויתי ואינם תלויים ב-$m$.

\section{הקשר למטריצות סיבוב \LR{D-Matrices}}
חיבור תנע זוויתי שקול מבחינת תורת החבורות לפירוק מכפלה של הצגות ($D^{(j)}$) לסכום ישר של הצגות בלתי פריקות (\LR{Clebsch-Gordan Series}).
המשוואה המרכזית המקשרת בין מכפלה של מטריצות סיבוב למקדמי \LR{CG} היא:
\[
D^{j_1}_{m_1 m'_1}(R) D^{j_2}_{m_2 m'_2}(R) = \sum_{J, M, M'} \braket{j_1 m_1 j_2 m_2}{J M} \braket{j_1 m'_1 j_2 m'_2}{J M'} D^{J}_{M M'}(R)
\]
משוואה זו מאפשרת לבטא מכפלה של פונקציות גל (שהן למעשה אלמנטים של מטריצות סיבוב, כמו ההרמוניות הספריות) כסכום לינארי של פונקציות גל עם תנע זוויתי כולל מוגדר.
